%!TeX root=../thesis.tex
\chapter{Introduction}
\label{ch:introduction}%

You can write an introduction here!

\section{Results}
We proof that the \textsc{Traveling Salesmen Problem}(TSP) is NP-complete.
\begin{verbatim}
... that the \textsc{Traveling Salesmen Problem}(TSP) is ...
\end{verbatim}
You can use abbreviations for your problem names to avoid redundancy.
Try to use existing abbreviations from related work if possible.



\section{Related Work}
% What has been done in existing literature?
% This is needed to provide context of the area you are working in.

% You could for example cite~\cite{limper2013pop} with
% \begin{verbatim}
% ... example cite~\cite{limper2013pop} with
% \end{verbatim}
% The $\sim$ keeps two word in the same line but also separates with a space.
% Note that there is no space before the $\sim$.

% \section{Hints}
% For titles see \url{https://en.wikipedia.org/wiki/Title_case}.

% Leave an empty line to start a new paragraph.
% As a general rule: Never force line breaks.
% \begin{verbatim}
% Never use \\ or \newline
% \end{verbatim}

% You can add and refer to figures like \cref{fig:logo}.
% \begin{verbatim}
% ... figures like \cref{fig:logo}.
% \end{verbatim}
% The \verb|\cref{}| command automatically detects the reference type.

% \begin{figure}
% 	\centering
% 	
\includegraphics[]{tu-bs_logo}
% 	\caption[Short description]{Long description: Notice that the List of Figures has a shorter description than this text right here. Isn't this fascinating?}
% 	\label{fig:logo}
% \end{figure}

% Generally, we prefer a graph $G$ is \emph{complete} over \textbf{complete} if \ldots
% \begin{verbatim}
% ... graph $G$ is \emph{complete} over \textbf{complete} if \ldots.
% \end{verbatim}

Das ist meine Übersicht zur Related work

\subsection{Delaunay Triangulations}
Als Grundlage für den Algorithms können Delaunay Triangulation betrachtet werden.
Diese Liefern eine Triangulation welche möglichst Gleichschenklige Dreiecke nutzt.
In dem Buch von de Berg et.al.~\cite{Delaunay_Triangulation} wird diese Variante der Triangulation genauer betrachtet.
In diesem Fall wird mit den Delaunay Triangulation das Problem gelöst, das höhendaten in einem zwei Dimensionalen Raum dargestellt werden.

Triangulation eignen sich aufgrund ihres Aufbaues gut um Höhendaten im zwei Dimensionalen Raum darzustellen (siehe Bild ~\ref{fig:terrain}).
Hier kann genutzt werden das Dreiecke verschiedene Innenwinkel und Kantenlängen haben können.
\begin{figure}[tb]
	\caption{\emph{A perspective view of a terrain} aus ~\cite[p.~183]{Delaunay_Triangulation}}
	\centering
	\label{fig:terrain}
\includegraphics[scale=0.2]{Delauney_Tarrain.png}
\end{figure}
Das Grundproblem ist hier das die höhendaten nur an bestimmten Punkten bekannt sind und somit die Restlichen Höhendaten erraten werden müssen.

\begin{itemize}
	\item für die darstellung Höhen raten
	\item Gleichschenklich sieht am besten aus
	\item ziel minimalen Winkel maximieren
	\item Formel für Kanten und Dreiecke bekommen
	\item drei Algos
\end{itemize}

\subsection{Flips in Triangulations}
Flips sind Möglichkeiten, um Kanten in einer Triangulation auszutauschen.
Da Triangulationen für die gleiche Knotenmenge, wie in~\cite{Delaunay_Triangulation} erwähnt, stets die gleiche Anzahl an Kanten besitzen, können keine Kanten hinzugefügt oder entfernt werden, um den Grad einer Triangulation zu verändern.
Daraus folgt, dass für jede entfernte Kante eine andere Kante eingefügt werden muss.
Eine besondere Art dieser Umstrukturierung von Kanten sind sogenannte Flips.
Hurtado et al.~\cite{Flipping_edges_in_triangulations} zeigten, dass jede Triangulation mindestens
\[\frac{n-4}{2}\]
flipbare Kanten besitzt, wobei \( n \) die Anzahl der Knoten bezeichnet.
Laut dem Paper \emph{Flips in planar graphs}~\cite{Flips_in_planar_graphs} lässt sich jede mögliche Triangulation in jede beliebige andere Triangulation überführen, indem die notwendigen Flips durchgeführt werden.
Daraus folgt, dass, wenn eine \emph{degree-constrained triangulation} möglich ist, diese mit den nötigen Flips konstruiert werden kann.

Eine weitere Art von Flips sind sogenannte \( k \)-Flips.
Diese stellen eine zusätzliche Möglichkeit dar, Kanten innerhalb einer Triangulation zu verschieben.
Dabei werden \( k \) Kanten entfernt und durch \( k \) neue Kanten ersetzt.
Normale Flips entsprechen somit \( 1 \)-Flips.
Der Vorteil von \( k \)-Flips liegt darin, dass mit einem Schritt größere strukturelle Veränderungen möglich sind.

\textbf{Absatz für Flips mit Binärbäumen hinzufügen, falls verwendet.}

\begin{figure}[ht]
  \centering
  \includegraphics[width=\linewidth]{drei_flips.png}
  \caption{\emph{Examples of edge move, edge rotation, and edge slide} aus~\cite[p.~70]{Flips_in_planar_graphs}}
  \label{fig:drei_flips}
\end{figure}

Drei beispielhafte Flips sind in Abbildung~\ref{fig:drei_flips} dargestellt.
Zum einen der \emph{edge move}, hierbei handelt es sich um das einfache Entfernen und anschließende Einfügen einer Kante.
Etwas stärker eingeschränkt ist die \emph{edge rotation}, bei der ein Endpunkt der Kante unverändert bleibt.
Die Rotation erfolgt also nur um einen Knoten.
Die dritte Variante, der \emph{edge slide}, schränkt die \emph{edge rotation} weiter ein.
Hier darf der veränderte Knoten nur ein Nachbar des ursprünglichen Knotens sein.

\subsection{degree constrained minimum spanning tree}
das sind viele infos

zum beispiel das brand and bound genutzt wird



\section{Overview}
If needed, write down where to find everything.
